\section{Answer reduction}\label{sec:answer-reduction}

In this section, we carry out the answer reduction.
Our main result will be to take the $\poly(n)$ question complexity, $O(2^n)$ answer complexity $\MIP^*$ protocol for $\succinctsquared$
given by \Cref{cor:succinct-sat-protocol-with-big-answer-size}
and convert it to one whose answer complexity is also $\poly(n)$;
this is \Cref{thm:main} below.

Our answer reduction will apply to any game with a value-$1$ real commuting EPR strategy.
We will require two properties of these strategies:
first, that they can be extended to strategies that pass subset tests with probability~$1$, as in \Cref{def:testing-a-code};
and second, that they are ``oracularizable".
We explain this second property in the next section. 

\subsection{Oracularization}

Our technique will not work for all entangled games but only for a
subset, for which a single prover can simulate both prover's actions
if required to. We call such games ``oracularizable'' games.


\begin{definition}
Given a two-player entangled game~$\game$,
its \emph{oracularization} is the game $C_{\mathrm{oracle}}(\game)$ given in \Cref{fig:oracle}.
If $\game$ is value-$1$,
then we call it \emph{oracularizable}, if $\val{C_{\mathrm{oracle}}(\game)} = 1$ as well. 
We also note that for \emph{any} game~$\game$, if $\val{\game} \leq 1 -
  \delta$, then $\val{C_{\mathrm{oracle}}(\game)} \leq 1 - O(\delta)$.
\end{definition}

{
  \floatstyle{boxed}
  \restylefloat{figure}
  \begin{figure}
    Given a game $\game$, sample a tuple $(\bx_0, \bx_1, \bC) \sim
    \game$, and flip two unbiased coins $\bb, \bc \sim \{0,1\}$.
    With probability $\tfrac{1}{2}$ each, perform one of the following two tests. 
    \begin{enumerate}
      \item \textbf{Verify:} Distribute the questions as follows:
      	\begin{itemize}
	\item[$\circ$] Player~$\bb$: send the pair $(\bx_0, \bx_1)$ and receive answers $(\ba_0, \ba_1)$.
	\item[$\circ$] Player~$\overline{\bb}$: send $\bx_{\bc}$ and receive an answer $\ba_2$.
	\end{itemize}
	\item \textbf{Consistency:} Play the consistency game with question $\bx_0, \bx_1$.
    \end{enumerate}
    Accept if $\ba_2 = \ba_{\bc}$ and $V(\bx_0, \bx_1, \ba_0, \ba_1) = 1$.
    \caption{The oracularized game $C_{\mathrm{oracle}}(\game)$.\label{fig:oracle}}
  \end{figure}
  }
  
\noindent
A real commuting EPR strategy allows ``Player~$\bb$" to sample both questions~$\bx_0$ and~$\bx_1$ simultaneously.
As a result, if a game~$\game$ has a value-$1$ real commuting EPR strategy, then it is oracularizable.

The value of oracularization is that when the verifier checks $V(\bx_0, \bx_1, \ba_0, \ba_1) = 1$,
both~$\ba_0$ and~$\ba_1$ come from the same prover rather than two different provers.
This seems like a minor change, but in fact it makes all the difference.
Our goal is to reduce the verifier's runtime by having the provers encode their answers using PCP technology.
When the answers come from \emph{both} provers,
the relevant piece of PCP technology is a \emph{distributed} PCP,
but it is known by a simple argument of Reingold that distributed PCPs do not exist
(see the discussion in~\cite{ARW17}).
The key difficulty comes from the fact that Alice needs to prepare her PCP proof without
knowing Bob's question and answer, and vice versa, and this turns out to be impossible in general.
On the other hand,
when the answers come from a single prover,
we can use traditional PCPs to implement the answer reduction, of which we have a variety of constructions.
We note that oracularized games \emph{do} still have checks between players,
but these are equality checks and will be easy to implement in the answer reduction regime.

\subsection{Probabilistically checkable proofs of proximity}

In this section, we introduce the main PCP technology we will use for our answer reduction.
In the oracularized game,
the provers want to convince us not just that $V(\cdot, \cdot, \cdot, \cdot)$
is satisfiable---which we already know to be true by construction---but that $(\bx_0, \bx_1, \ba_0, \ba_1)$
is a particular assignment which satisfies it.
For this, we need a stronger notion of a PCP called a \emph{probabilistically checkable proof of proximity (PCPP)}.
These allow one to check that an input~$x$ is close to a satisfying assignment of a circuit~$C$ (hence the ``proximity")
by making a small number of queries to~$x$.
These were originally introduced in the independent works of~\cite{BGH+06}
and~\cite{DR06} (where they were called \emph{assignment testers}).

In our case, we will need even stronger PCPPs in which the verifier is not only query-efficient but \emph{time}-efficient as well.
The history of these time-efficient PCPPs goes back to the original proof of $\MIP = \nexp$ and the various attempts to ``scale it down"~\cite{OD05}.
The most famous line of research considered proof systems in which the verifier's query complexity is restricted,
and this eventually led to the proof of the PCP theorem~\cite{AS98,ALM+98}.
A parallel line of research considered proof systems in which the verifier's runtime is restricted (so-called ``transparent" proofs)~\cite{BFLS91}.
The latter of these was revisited in the work of Ben-Sasson et al.~\cite{BGH+05},
who showed that both lines of research could be remerged in the ``scaled down" setting
by constructing a PCPP in which the verifier is both query-efficient \emph{and} time-efficient.
Though their main result is actually sufficient for our purposes,
we will cite the work of Mie~\cite{Mie09}, which improves on their result in the regime we care about.
Finally, we note the work of Meir~\cite{Mei14}, who reproves the bounds of Ben-Sasson et al.~\cite{BGH+05} using combinatorial methods.

To our knowledge, ours is the first use of a time-efficient PCPP specifically for its time-efficient properties in the quantum literature.
Natarajan and Vidick~\cite{NV18a} used the time-efficient PCPP of~\cite{BGH+05} to prove the quantum games PCP conjecture,
but the property they needed was not that it was time-efficient,
but that the bits of the proof are linear functions of the bits of the assignment.
We note that we do not need this property here.

In this literature, it is common to consider ``pair languages" consisting of strings $(x, y)$
in which~$x$ is small and given to the verifier and~$y$ is large and accessible only through query access.
This maps perfectly onto our scenario,
in which the verifier supplies the ``small" questions $\bx_0, \bx_1$
and the prover supplies the ``large" answers $\ba_0, \ba_1$.

\begin{definition}
A \emph{pair language} $L$ is a subset of $\{0, 1\}^* \times \{0, 1\}^*$.
Given $x \in \{0, 1\}^*$, we write $L_x = \{y \in \{0, 1\}^* \mid (x,y) \in L\}$.
\end{definition}

The next two definitions state the notion of an efficient PCPP verifier.

\begin{definition}[{\cite[Definition 2.1]{BGH+05}}]
Let $r, q:\Z^+ \rightarrow \Z^+$ and $t:\Z^+ \times \Z^+ \rightarrow \Z^+$.
An \emph{$(r, q, t)$-restricted PCPP verifier} is a probabilistic machine that, given a string~$x$ (called the \emph{explicit input})
and a number~$K$ (in binary) as well as oracle access to an \emph{implicit input} $y \in \{0, 1\}^{K}$ and to a \emph{proof oracle} $\pi \in \{0, 1\}^*$,
tosses $r(|x|+K)$ coins, queries the oracles $(y, \pi)$ for a total of $q(|x|+K)$ symbols,
runs in time $t(|x|, K)$, and outputs a Boolean verdict in $\{\mathrm{accept}, \mathrm{reject}\}$.
\end{definition}

\begin{definition}[{\cite[Definition 2.2]{BGH+05}}]
For functions $r, q:\Z^+ \rightarrow \Z^+$, $t:\Z^+\times \Z^+ \rightarrow \Z^+$,
and constants
$s, \gamma \in  [0, 1]$,
a pair language $L \subseteq \{0, 1\}^* \times \{0, 1\}^*$ is in $\mathrm{PCPP}_{s, \gamma}[r, q, t]$
if there exists an $(r, q, t)$-restricted PCPP verifier~$V$ with the following properties:
\begin{itemize}
\item[$\circ$] \textbf{Completeness:} If $(x, y) \in L$ then there exists a $\pi$ such that $\Pr_R[\text{$V^{y, \pi}(x, |y|; R)$ accepts}] = 1$,
	where $V^{y, \pi}(x, |y|; R)$ denotes the decision $V$ on input $(x, |y|)$,
	oracle access to $(y, \pi)$, and coin tosses $R$.
\item[$\circ$] \textbf{Soundness:} If $(x, y)$ is such that $y$ is $\gamma$-far from $L_x \cap \Sigma^{|y|}$,
	then for every $\pi$ it holds that $\Pr_R[\text{$V^{y, \pi}(x, |y|;R)$ accepts}] \leq s$.
\end{itemize}
\end{definition}

Mie's time-efficient PCPP is states as follows.

\begin{theorem}[{\cite[Theorem 1]{Mie09}}]\label{thm:pcpp-mie}
Suppose that $L$ is a pair language in $\ntime(T)$
for some non-decreasing function $T : \Z^+ \rightarrow \Z^+$.
Then, for every two constants $s, \gamma >0$,
we have $L \in \mathrm{PCPP}_{s, \gamma}[r, q, t]$, for
\begin{itemize}
\item[$\circ$] Randomness complexity $r(m) = \log_2 T(m) + O(\log\log T(m))$.
\item[$\circ$] Query complexity $q(m) = O(1)$,
\item[$\circ$] Verification time $t(n, K) = \poly(n, \log K, \log T(n + K))$.
\end{itemize}
\end{theorem}

We note that this is in fact a much stronger than what we will actually need.
In particular, we will only apply this to languages $L$ in \emph{deterministic} $\mathsf{TIME}(T)$,
which are trivially in $\ntime(T)$.

\subsection{Composing with an error-correcting code}\label{sec:composing-with-an-error-correcting-code}

The verifier in a PCPP rejects any input which is $\gamma$-far from an accepting input,
but of course we want our verifier to reject \emph{all} non-accepting inputs, no matter their distance.
To do this, we will (i) encode the verifier's inputs using an error-correcting code
and (ii) check that the inputs are properly encoded (using, for example, the low-degree test).
This approach of composing a PCPP with an error-correcting code is standard
and stretches back in spirit to the transparent proofs of~\cite{BFLS91} (see the discussion of this in~\cite{BGH+06}).

Now we show how to compose an $\MIP^*$ game with an error-correcting code.

\begin{definition}[Error-correcting the provers' answers]\label{def:error-correcting-answers}
Let $V = (\mathrm{Alg}_{\mathrm{Q}}, \mathrm{Alg}_{\mathrm{A}})$ be an $\MIP^*$ verifier (the language it verifies is not important).
Suppose on inputs of size~$n$ it has question length~$\ell_{\mathrm{Q}}(n)$ answer length~$\ell_{\mathrm{A}}(n)$.
Write $L_{\mathrm{A}}$ for the language decided by $\mathrm{Alg}_{\mathrm{A}}$.
Let $\codec_k = (\codee_k,\coded_k, \codes_k)$ be a $(k, m, q, \eta, t_{\mathrm{Dec}}, t_{\mathrm{Emb}})$-efficient code family with decoding algorithm $\mathrm{Alg}_{\mathrm{Dec}}$.
Then $L_{\mathrm{A}} \circ \codec$ is a new language defined as follows:
suppose $(\mathsf{input}, x_0, x_1, y_0, y_1) \in L_{\mathrm{A}}$.
Let $n$ be the length of $\mathsf{input}$ and $\ell = \ell_{\mathrm{A}}(n)$.
Then $(\mathsf{input}, x_0, x_1, \codee_{\ell}(y_0), \codee_{\ell}(y_1)) \in L_{\mathrm{A}} \circ \mathsf{Code}$.
\end{definition}

Now, we prove a couple of properties about the composed verifier.
First, we show that its runtime is not much slower than the original verifier's.

\begin{proposition}[Runtime of the composed verifier]\label{prop:composed-runtime}
Let~$V$ and~$\codec_k$ be as in \Cref{def:error-correcting-answers}.
Suppose $\mathrm{Alg}_{\mathrm{A}}$ runs in time $T(n)$.
Then there is an algorithm, which we denote $\mathrm{Alg}_{\mathrm{A}} \circ \mathsf{Code}$,
deciding the language $L_{\mathrm{A}} \circ \mathsf{Code}$.
In addition, on inputs $(\mathsf{input}, x_0, x_1, z_0, z_1)$
in which $|\mathsf{input}| = n$,
$|x_0| = |x_1| = \ell_{\mathrm{Q}}(n)$,
and $|z_0| = |z_1| = m(\ell_{\mathrm{A}}(n))$,
the algorithm runs
in time $T(n) + t_{\mathrm{Dec}}(\ell_{\mathrm{A}}(n))$.
\end{proposition}

\begin{proof}
On input $(\mathsf{input}, x_0, x_1, z_0, z_1)$, we define the action of $\mathrm{Alg}_{\mathrm{A}} \circ \mathsf{Code}$ as follows.
\begin{enumerate}
\item Compute $n$, the length of $\mathsf{input}$. Set $\ell := \ell_{\mathrm{A}}(n)$.
\item Check that~$z_0$ and~$z_1$ have length~$m(\ell)$. If they don't, reject.
\item Compute $y_0 = \mathrm{Alg}_{\coded}(\ell, z_0)$ and $y_1 = \mathrm{Alg}_{\coded}(\ell, z_1)$. If either $y_0$ or $y_1$ is $\bot$, reject.\label{item:decode-that-string}
\item Otherwise, we know that $y_0, y_1 \in \{0, 1\}^\ell$.  Run $\mathrm{Alg}_{\mathrm{A}}(\mathsf{input}, x_0, x_1, y_0, y_1)$. Accept if it accepts, and reject if it rejects.\label{item:run-that-algorithm}
\end{enumerate}
It is immediate that $\mathrm{Alg}_{\mathrm{A}} \circ \mathsf{Code}$ computes $L_{\mathrm{A}} \circ \mathsf{Code}$.
As for the time complexity, \Cref{item:decode-that-string} runs in time $t_{\mathrm{Dec}}(\ell_{\mathrm{A}}(n))$ 
and \Cref{item:run-that-algorithm} runs in time $T(n)$.
Combined, these two give the bound in the proposition statement.
\end{proof}

Next, we show that this construction solves the ``problem" discussed at the beginning of the section,
namely that if we perform answer reduction by replacing $\mathrm{Alg}_{\mathrm{A}} \circ \mathsf{Code}$ with a PCPP verifier,
rather than just $\mathrm{Alg}_{\mathrm{A}}$, then the verifier will reject \emph{all} inputs which are not in the language,
not just those which are $\delta$-far, provided that those inputs are encoded as per \Cref{def:error-correcting-answers}.

\begin{proposition}\label{prop:new-soundness}
Let~$V$ and~$\codec_k$ be as in \Cref{def:error-correcting-answers}.
Let $s, \gamma > 0$ be constants,
and let $V_{\mathrm{PCPP}}$ be the PCPP verifier for the language $L_{\mathrm{A}} \circ \mathsf{Code}$
guaranteed by \Cref{thm:pcpp-mie} with these parameters.
Suppose that $1-\eta(k) \geq 2\gamma$ for all~$k$.
Then we have the following soundness condition.
\begin{itemize}
\item[$\circ$] \textbf{Soundness:} Consider $(\mathsf{input}, x_0, x_1, z_0, z_1)$ for $\mathsf{input}$ of length~$n$, $x_0$ and~$x_1$ of length~$\ell_{\mathrm{Q}}(n)$,
	and $z_0, z_1 \in \codes_{\ell}$, for $\ell := \ell_{\mathrm{A}}(n)$.
	Suppose this does not correspond to the encoding of an accepting assignment in $L_{\mathrm{A}}$.
	In other words, suppose that there are no $y_0, y_1 \in \{0, 1\}^\ell$ such that $(\mathsf{input}, x_0, x_1, y_0, y_1)$ is in $L_{\mathrm{A}}$
	and $z_0 = \codee_{\ell}(y_0), z_1 = \codee_{\ell}(y_1)$.
	Then $V_{\mathrm{PCPP}}$ accepts $(\mathsf{input}, x_0, x_1, z_0, z_1)$ with probability at most~$s$.
	In math, for every $\pi$ it holds that
	\begin{equation*}
		\Pr_R[\text{$V_{\mathrm{PCPP}}^{z_0, z_1, \pi}(\mathsf{input}, x_0, x_1, |z_0| + |z_1|; R)$ accepts}] \leq s.
	\end{equation*}
\end{itemize}
\end{proposition}
\begin{proof}
Given $(\mathsf{input}, x_0, x_1, z_0, z_1)$, write $A:= (L_{\mathrm{A}} \circ \mathsf{Code})_{\mathsf{input}, x_0, x_1} \cap \F_{q(\ell)}^{|z_0| + |z_1|}$.
By assumption, $(z_0, z_1)$ is not in~$A$.
Using this, we would like to show that $(z_0, z_1)$ is in fact $\gamma$-far from~$A$, in which case the PCPP verifier accepts with probability at most~$s$.

To do this, suppose $(z_0', z_1') \in A$.
By design, there exists $y_0', y_1' \in \{0, 1\}^\ell$ such that $z_0' = \codee(y_0')$ and $z_1' = \codee(y_1')$.
This means that $z_0', z_1' \in \codes_{\ell}$.
On the other hand, since~$(z_0, z_1)$ is not in~$A$, we must have either $z_0' \neq z_0$ or $z_1' \neq z_1$ (or both).
We will assume the first without loss of generality.
Then by the distance property of the code, since $z_0, z_0' \in \codes_{\ell}$,
their normalized Hamming distance is at least $1-\eta(\ell) \geq 2\gamma$.
This immediately means that $(z_0,z_1)$ and $(z_0', z_1')$.
are at least $\gamma$-far apart, and we are done.
\end{proof}


\subsection{The answer reduction protocol}

We are almost ready to state the answer reduction protocol.
Before doing so, we discuss one final nuisance, which is that we will also need the prover to encode their \emph{proof} with an error-correcting code. 
The reason is that we would like to query the proof on a view~$\bJ$ sampled by the PCPP verifier.
However, the prover might cheat and respond based only on the view~$\bJ$ rather than a global proof~$\pi$.
To prevent this, we force them to commit to a global error-correcting encoding of their proof~$\pi$ using a tester as in \Cref{def:testing-a-code}.
Then, we use the fact that the error-correcting code \emph{embeds} their string to allow us to extract the view~$\bJ$ by asking for the coordinates in~$\mu(\bJ)$.

We now state the answer reduction protocol.

\begin{definition}\label{def:answer-reduction-game}
We instantiate the answer-reduced $\mip^*$ protocol with the following algorithms and parameters.
\begin{itemize}
\item[$\circ$] Let $V = (\mathrm{Alg}_{\mathrm{Q}}, \mathrm{Alg}_{\mathrm{A}})$ be an $\MIP^*$ verifier for a language~$L$.
		Write $L_{\mathrm{A}}$ for the language decided by $\mathrm{Alg}_{\mathrm{A}}$.
		Suppose on inputs of size~$n$, the verifier~$V$ has question length~$\ell_{V,\mathrm{Q}}(n)$, answer length~$\ell_{V,\mathrm{A}}(n)$,
		question time~$t_{V,\mathrm{Q}}(n)$, and answer time~$t_{V,\mathrm{A}}(n)$.
\item[$\circ$] Let $\codec_k = (\codee_k,\coded_k, \codes_k)$ be a $(k, m, q, \eta, t_{\mathrm{Dec}}, t_{\mathrm{Emb}})$-efficient code family
		with decoding algorithm $\mathrm{Alg}_{\mathrm{Dec}}$ and embedding $\mu_k$.
\item[$\circ$] Let $\game_k$ be a game which tests for $\codec_k$ with robustness $\chi_k(\eps)$.
		Suppose it has question length~$\ell_{\game, \mathrm{Q}}(k)$, answer length~$\ell_{\game,\mathrm{A}}(k)$,
		question time~$t_{\game,\mathrm{Q}}(k)$, and answer time~$t_{\game,\mathrm{A}}(k)$.
\item[$\circ$] Let $s, \delta > 0$ be constants, and let $V_{\mathrm{PCPP}}$ be the PCPP verifier for the language $L_{\mathrm{A}} \circ \mathsf{Code}$
			guaranteed by \Cref{thm:pcpp-mie} with these parameters.
			Suppose on inputs of size~$n$ it has proof length~$\ell_\pi(n)$.
			By \Cref{prop:composed-runtime}, $L_{\mathrm{A}} \circ \codec$ is in time
			$t_{\mathrm{compose}}(n) = t_{V, \mathrm{A}}(n) + t_{\mathrm{Dec}}(\ell_{V,\mathrm{A}}(n))$.
			We can therefore write $V_{\mathrm{PCPP}}$'s verification time as
			\begin{equation*}
				t_{\mathrm{PCPP}}(n) = \poly(n + \ell_{V, \mathrm{Q}}(n), \log(m(\ell_{V, \mathrm{A}}(n))), \log(t_{\mathrm{compose}}(n))).
			\end{equation*}
			Finally, $\ell_\pi(n) = t_{\mathrm{compose}}(n) \cdot \mathrm{poly log}(t_{\mathrm{compose}}(n))$.
\end{itemize}
Write $\ell_1 := \ell_{V, \mathrm{A}}(n)$ and $\ell_2 := \ell_\pi(n)$.
Then the \emph{answer reduction game} $\game_{\mathrm{answer}}(\mathsf{input};V, \codec, \game, s, \delta)$
is given in \Cref{fig:answer-reduction}.
We write $V_{\mathrm{answer}}$ for the corresponding verifier.
\end{definition}

{
	\floatstyle{boxed}
	\restylefloat{figure}
	\begin{figure}Flip two unbiased coins $\bb, \bc \sim \{0, 1\}$. 
    		Sample questions $(\bx_0, \bx_1) \sim \mathrm{Alg}_{\mathrm{Q}}(\mathsf{input})$.
		Sample a view $\bI_0, \bI_1, \bJ \sim V_{\mathrm{PCPP}}(\mathsf{input}, \bx_0, \bx_1)$.
		Set $\bJ' = \mu_{\ell_2}(\bJ)$.
		Select $\bi_0, \bi_1 \in [m(\ell_1)]$ and $\bj \in [m(\ell_\pi(n))]$ uniformly at random.
		Set $\bT_0 = \bI_0 \cup \{\bi_0\}$, $\bT_1 = \bI_1 \cup \{\bi_1\}$, and $\bU = \bJ' \cup \{\bj\}$.
	With probability $\frac{1}{8}$ each, perform one of the following eight tests.
	\begin{enumerate}
	\item \textbf{Verify:}
		Distribute the question as follows:
			\begin{itemize}
			\item[$\circ$] Player~$\bb$: give $(\bx_0, \bx_1)$, $\bT_0, \bT_1, \bU$; receive $\ba_0, \ba_1, \ba_2$.
			\end{itemize}
		Accept if $V_{\mathrm{PCPP}}(\mathsf{instance}, \bx_0, \bx_1)$ accepts on $\ba_0|_{\bI_0}, \ba_1|_{\bI_1}, \ba_2|_{\bJ'}$.
	\item \textbf{Cross checks:}
	\begin{enumerate}
	\item \textbf{Consistency test:} Distribute the questions as follows:
			\begin{itemize}
			\item[$\circ$] Player~$\bb$: give $(\bx_0, \bx_1)$, $\bT_0, \bT_1, \bU$; receive $\ba_0, \ba_1, \ba_2$.
			\item[$\circ$] Player~$\overline{\bb}$: give $(\bx_0, \bx_1)$, $\bT_0, \bT_1, \bU$; receive $\ba_0', \ba_1', \ba_2'$.
			\end{itemize}
		Accept if $\ba_0 = \ba_0'$, $\ba_1 = \ba_1'$, and $\ba_2 = \ba_2'$.
	\item \textbf{Answer cross-check:} Distribute the questions as follows:
			\begin{itemize}
			\item[$\circ$] Player~$\bb$: give $(\bx_0, \bx_1)$, $\bT_0, \bT_1, \bU$; receive $\ba_0, \ba_1, \ba_2$.
			\item[$\circ$] Player~$\overline{\bb}$: give $\bx_{\bc}, \bT_{\bc}'$; receive $\ba_{\bc}'$.
			\end{itemize}
		Accept if $\ba_{\bc} = \ba_{\bc}'$.
	\item \textbf{Proof cross-check:} Distribute the questions as follows:
			\begin{itemize}
			\item[$\circ$] Player~$\bb$: give $(\bx_0, \bx_1)$, $\bT_0, \bT_1, \bU$; receive $\ba_0, \ba_1, \ba_2$.
			\item[$\circ$] Player~$\overline{\bb}$: give $\bx_0, \bx_1, \bU$; receive $\ba_{2}'$.
			\end{itemize}
		Accept if $\ba_2 = \ba_2'$.
	\end{enumerate}
	\item \textbf{Code checks:}
	\begin{enumerate}
	\item \textbf{Answer code check:}
		Sample questions $(\bw_0, \bw_1) \sim \game_{\ell_1}(\bT_{\bc})$.
		Distribute the questions as follows:
			\begin{itemize}
			\item[$\circ$] Player~$\bb$: give $\bx_{\bc}, \bw_{0}$; receive $\ba_0$.
			\item[$\circ$] Player~$\overline{\bb}$: give $\bx_{\bc}, \bw_{1}$; receive $\ba_1$.
			\end{itemize}
		Accept if $\game_{\ell_1}(\bT_{\bc})$ accepts on $\ba_0, \ba_1$.
	\item \textbf{Proof code check:}
		Sample questions $(\bw_0, \bw_1) \sim \game_{\ell_2}(\bU)$.
		Distribute the questions as follows:
			\begin{itemize}
			\item[$\circ$] Player~$\bb$: give $\bx_0, \bx_1$, $\bw_{0}$; receive $\ba_0$.
			\item[$\circ$] Player~$\overline{\bb}$: give $\bx_0, \bx_1$, $\bw_{1}$; receive $\ba_1$.
			\end{itemize}
		Accept if $\game_{\ell_2}(\bU)$ accepts on $\ba_0, \ba_1$.
	\end{enumerate}
    \end{enumerate}
    \caption{The answer reduction game $\game_{\mathrm{answer}}(\mathsf{input};V, \codec, \game, s, \delta)$.\label{fig:answer-reduction}}
  \end{figure}
}

\begin{theorem}\label{thm:answer-reduction}
Suppose~$V$, $\codec$, $\game$, and $V_{\mathrm{PCPP}}$ are as in \Cref{def:answer-reduction-game}.
Suppose $s, \gamma$ are chosen to be constants such that $\eta(k) \geq 2\gamma$ for all~$k$.
Suppose further that~$V$ has the following property: for any input in~$L$,
the provers have a real commuting EPR strategy with value~$1$.
Then~$V_{\mathrm{answer}}$ is also an $\MIP^*$ verifier for~$L$ with the following two conditions:
\begin{itemize}
\item[$\circ$] (Completeness) If $\mathsf{input} \in L$, then there is a value-$1$ strategy.
\item[$\circ$] (Soundness) Given $\mathsf{input}$, suppose there is a strategy with value $1-\eps$.
			Then there is a strategy for~$V$ on input $\mathsf{input}$ with value $1-\delta(\eps)$, where $\delta(\eps)$ is given by
\begin{equation*}
\delta(\eps) := \poly(\chi_{\ell_1}(\poly(\eps)), \chi_{\ell_2}(\poly(\eps)), \eta(\ell_1), \eta(\ell_2)).
\end{equation*}
\end{itemize}
Hence, if we choose our parameters so that $1-\delta(\eps)$ is greater than the soundness of~$V$,
this implies that $V_{\mathrm{answer}}$ is an $\MIP^*$ verifier for~$L$ with soundness $1-\eps$.

Furthermore, the question and answer lengths and runtimes are dominated by two subroutines: the ``Verify" subroutine $S_1$ 
and the ``Code Check" subroutine $S_2$ (consisting of both the answer code check and the proof code check).
The complexity of the Verify subroutine is
\begin{align*}
\qlength{S_1} &= O(\ell_{V, \mathrm{Q}}(n) + \log(m(\ell_{V,\mathrm{A}}(n)))+ \log(m(\ell_{\pi}(n)))),\\
\alength{S_1} &= O(\log(q(\ell_{V,\mathrm{A}}(n))) + \log(q(\ell_\pi(n)))),\\
\qtime{S_1} &= O\left(t_{V,\mathrm{Q}}(n) + t_{\mathrm{PCPP}}(n) + t_{\mathrm{Emb}}(\ell_{\pi}(n))\right),\\
\atime{S_1} &= O(t_{\mathrm{PCPP}}(n)).
\end{align*}
In addition, the complexity of the Code Check subroutine is
\begin{align*}
\qlength{S_2} &= O(\ell_{\game, \mathrm{Q}}(\ell_{V, \mathrm{A}}(n)) + \ell_{\game, \mathrm{Q}}(\ell_\pi(n)) + \ell_{V, \mathrm{Q}}(n)),\\
\alength{S_2} &= O(\ell_{\game, \mathrm{A}}(\ell_{V, \mathrm{A}}(n)) + \ell_{\game, \mathrm{A}}(\ell_\pi(n))),\\
\qtime{S_2} &=  O(t_{\game, \mathrm{Q}}(\ell_{V, \mathrm{A}}(n)) + t_{\game, \mathrm{Q}}(\ell_\pi(n)) + t_{V, \mathrm{Q}}(n) + t_{\mathrm{Emb}}(\ell_{\pi}(n))),\\
\atime{S_2} &= O(t_{\game, \mathrm{A}}(\ell_{V, \mathrm{A}}(n)) + t_{\game, \mathrm{A}}(\ell_\pi(n))).
\end{align*}
Thus, the complexity of the overall protocol is the sum of these two.
\end{theorem}
\begin{proof}
The fact  that~$S_1$ and~$S_2$ dominate the lengths and runtimes of the protocol
is because~$S_1$ dominates the lengths and runtimes of the two cross-check subroutines,
whose questions and answers are subsets of those in~$S_1$.
Now we compute the complexity of~$S_1$.
\begin{itemize}
\item[$\circ$] \textbf{Question length:} The pair $(\bx_0, \bx_1)$ has total length $\ell_{V,\mathrm{Q}}(n)$ by definition.
		The pair $\bI_0, \bI_1$ are subsets of indices of constant size into each of the implicit inputs of $L_{\mathrm{A}} \circ \codec$,
		which are supposed to be encodings of strings of size $\ell_{V, \mathrm{A}}(n)$.
		Hence, the encodings have size $m(\ell_{V, \mathrm{A}}(n))$, and so each input is specified with the log of this many bits.
		Finally, $\bJ$ is a constant-sized set of indices into a proof of size $\ell_{\pi}(n)$, and $\mu(\bJ)$ converts these into indices into an encoding of of this proof.
		As the encoding has length $m(\ell_{\pi}(n))$, each index can be specified with the log of this many bits.
\item[$\circ$] \textbf{Answer length:} The strings $\ba_0, \ba_1$ contains values from an error-correcting code with alphabet $q(\ell_{V, \mathrm{A}}(n))$,
		and the string~$\ba_2$ contains values from an error-correcting code with alphabet $q(\ell_{\pi}(n))$.
\item[$\circ$] \textbf{Question time:} The running time of $\mathrm{Alg}_{\mathrm{Q}}$ is $t_{V,\mathrm{Q}}(n)$.
		The running time of $V_{\mathrm{PCPP}}$ is $t_{\mathrm{PCPP}}(n)$.
		Finally, the running time to compute $\mu(\bJ)$ given~$\bJ$ is $t_{\mathrm{Emb}}(\ell_{\pi}(n))$.
\item[$\circ$] \textbf{Answer time:} The running time is simply the running time of $V_{\mathrm{PCPP}}$, i.e.\ $t_{\mathrm{PCPP}}(n)$.
\end{itemize}
As for the complexity of~$S_2$, it just performs the code tester~$\game_k$ for message lengths $k = \ell_{V, \mathrm{A}}(n)$ and $\ell_\pi(n)$ and so inherits the lengths and runtimes of the tester for these two values of~$k$, except on top of that it also has to sample~$(\bx_0, \bx_1)$ and compute~$\bJ'$.
Sampling~$(\bx_0, \bx_1)$ takes time~$t_{V,\mathrm{Q}}(n)$ and contributes $O(\ell_{V, \mathrm{Q}}(n))$ to the question lengths,
and computing~$\bJ'$ takes time~$t_{\mathrm{Emb}}(\ell_{\pi}(n))$.

\paragraph{Completeness.}
Suppose $\mathsf{input}$ is in~$L$.
Then there is a real commuting EPR strategy $(\psi, M)$ with value~$1$ for~$V$ on $\mathsf{input}$.
We will use this to demonstrate a value-$1$ strategy for $V_{\mathrm{answer}}$.
This will be the strategy $(\psi, G)$ which uses the same EPR state $\ket{\psi}$ and has measurement matrices~$G$ defined as follows.

Fix questions $x_0, x_1$, $T_0, T_1$, and $U$.  We begin by defining the simplest measurement,
\begin{equation}\label{eq:i-dont-know-what-to-call-this}
G^{x_c, T_c}_{a_c} := M^{x_c}_{[\codee_{\ell_1}(z_c)|_{T_c} = a_c]}.
\end{equation}
If Alice and Bob measure with $M^{x_0}$ and $M^{x_1}$ and receive strings $z_0, z_1$,
then because this strategy is value~$1$,
we will always have $V(\mathsf{input}, x_0, x_1, z_0, z_1) = 1$.
As a result, there always exists \emph{some} proof for $V_{\mathrm{PCPP}}$
that $(\mathsf{input}, x_0, x_1, \codee_{\ell_1}(z_0), \codee_{\ell_1}(z_1))$ is in $L_{\mathrm{A}} \circ \mathsf{Code}$.
We denote this proof $\pi(x_0, x_1, z_0, z_1)$; if there are multiple such proofs, we pick one arbitrarily.
Then we define the measurement
\begin{equation}\label{eq:i-dont-know-what-to-call-this-part-two-return-of-the-bat}
G^{x_0, x_1, U}_{a_2} := (M^{x_0} \cdot M^{x_1})_{[\codee_{\ell_2}(\pi(x_0, x_1, z_0, z_1))|_{U} = a_2}.
\end{equation}
Next, we define the measurement
\begin{equation*}
G^{x_0, x_1, T_0, T_1, U}_{a_0, a_1, a_2}
:= (M^{x_0}\cdot M^{x_1})_{[\codee_{\ell_1}(z_0)|_{T_0}, \codee_{\ell_1}(z_1)|_{T_1}, \codee_{\ell_2}(\pi(x_0, x_1, z_0, z_1))|_{U} = a_0, a_1, a_2]}.
\end{equation*}
Now, via \Cref{eq:i-dont-know-what-to-call-this,eq:i-dont-know-what-to-call-this-part-two-return-of-the-bat},
the~$G$ measurement is exactly of the form required by \Cref{def:testing-a-code}.
As a result, it can be extended to a measurement which passes the answer and proof code checks with probability~$1$.
Performing this extension concludes the design of the strategy.

By construction, this strategy passes the answer and proof code checks with probability~$1$.
As for the remaining tests, let us begin with the answer cross-check in the case of $\bc = 0$, the other case being symmetric.
Because~$M$ is a real commuting EPR strategy,
by \Cref{fact:heh-heh-heh-gonna-make-anand-prove-this-so-i-can-take-the-day-off}
we have that $M^x_a \otimes I_{\reg{Bob}} \consistency_0 I_{\reg{Alice}} \otimes M^x_a$ for any distribution on~$x$.
If we consider the measurement $(M^{x_0} \cdot M^{x_1})_{z_0, z_1}$, then $(M^{x_0} \cdot M^{x_1})_{z_0} = M^{x_0}_{z_0}$. As a result,
\begin{equation*}
(M^{x_0} \cdot M^{x_1})_{z_0} \otimes I_{\reg{Bob}}
\consistency_0 I_{\reg{Alice}} \otimes M^{x_0}_{z_0}.
\end{equation*}
Finally, by data processing (\Cref{fact:specialize-the-simeq}), this implies that
\begin{equation*}
(M^{x_0} \cdot M^{x_1})_{[\codee_{\ell_1}(z_0)|_{T_0}=a_0]} \otimes I_{\reg{Bob}}
\consistency_0 I_{\reg{Alice}} \otimes M^{x_0}_{[\codee_{\ell_1}(z_0')|_{T_0} = a_0]}.
\end{equation*}
But this is equivalent to saying that $G^{x_0, x_1, T_0, T_1, U}_{a_0} \otimes I_{\reg{Bob}} \consistency_0 I_{\reg{Alice}} \otimes G^{x_0, T_0}_{a_0}$,
which implies passing the cross-check test with probability~$1$.
A similar argument holds for the other tests, with the exception of the verification step.

Consider the measurement $M^{x_0}_{z_0} \cdot M^{x_1}_{z_1}$.
By construction and the properties of the PCPP verifier, if this measurement always outputs~$z_0, z_1$ such that $V(\mathsf{input}, x_0, x_1, z_0, z_1) = 1$,
then the~$G$ strategy always passes the verify step.
But because~$M$ is a real commuting EPR strategy,
$M^x_z \otimes I_{\reg{Bob}} \consistency_0 I_{\reg{Alice}} \otimes M^x_z$,
which implies that
\begin{equation*}
M^{x_0}_{z_0} \otimes M^{x_1}_{z_1}
\approx_0 M^{x_0}_{z_0} \cdot M^{x_1}_{z_1} \otimes I_{\reg{Bob}}.
\end{equation*}
Thus, these two measurements have the same output distribution.
But the left-hand side always outputs~$z_0, z_1$ which satisfy the verifier, because this strategy passes the verifier with probability~$1$.
This concludes the completeness step.


\paragraph{Soundness.}
Suppose $\mathsf{input}$ is not in~$L$.
Let $(\psi, M)$ be a strategy that passes with probability $1-\eps$.


\paragraph{Code checks.}
Passing the overall test with probability $1-\eps$
means the strategy passes the answer code check with probability $1-8\eps$.
Given values $c $, $x_{c}$, write $1-\eps_{c, x_{c}}$ for the probability the code check passes conditioned on these values.
Then with probability at least $1-8\eps^{1/2}$, $\eps_{c, x_{c}} \leq \eps^{1/2}$.
When this occurs, \Cref{thm:low-degree-code-tester} implies that there exists a measurement $\{G^{x_c}_w\}_w$ with outcomes in $\codes_{\ell_1}$ such that
\begin{equation*}
M^{x_c, T_c}_a \otimes I_{\reg{Bob}}\consistency_{\delta(\eps)} I_{\reg{Alice}} \otimes G^{x_c}_{[w|_{T_c} = a]}
\end{equation*}
with respect to the distribution of~$\bT_c$ conditioned on~$c$ and~${x}_{c}$.
When this does \emph{not} occur, we still can assume such a measurement so that
\begin{equation*}
M^{x_c, T_c}_a \otimes I_{\reg{Bob}}\consistency_{1} I_{\reg{Alice}} \otimes G^{x_c}_{[w|_{T_c} = a]}
\end{equation*}
trivially, by \Cref{fact:trivial-upper-bound-approx-delta}. Thus, if we average over $\bc$ and $\bx_{\bc}$,
\begin{equation}\label{eq:almost-at-tweedledee}
M^{x_c, T_c}_a \otimes I_{\reg{Bob}}\consistency_{\delta(\eps)} I_{\reg{Alice}} \otimes G^{x_c}_{[w|_{T_c} = a]}
\end{equation}
with respect to the distribution on $\bc, \bx_{\bc}, \bT_{\bc}$.
A similar argument with respect to the consistency guarantee of \Cref{thm:low-degree-code-tester} implies that
\begin{equation}\label{eq:consistency-of-the-encoding}
G^{x_c}_w \otimes I_{\reg{Bob}} \consistency_{\delta(\eps)} I_{\reg{Alice}} \otimes G^{x_c}_w.
\end{equation}
By \Cref{fact:specialize-the-simeq}, this implies that
\begin{equation*}
G^{x_c}_{[w|_{T_c} = a]} \otimes I_{\reg{Bob}} \consistency I_{\reg{Alice}} \otimes G^{x_c}_{[w|_{T_c} = a]}.
\end{equation*}
As a result, if we apply \Cref{fact:agreement} to this and \Cref{eq:almost-at-tweedledee} and then use the triangle inequality (\Cref{fact:triangle}), we conclude
\begin{equation}\label{eq:tweedledee}
M^{x_c, I_c}_a \otimes I_{\reg{Bob}}\approx_{\delta(\eps)}G^{x_c}_{[w|_{I_c} = a]} \otimes  I_{\reg{Bob}}
\end{equation}
with respect to the distribution on $\bc, \bx_{\bc}, \bI_{\bc}$.

Applying a similar argument yet again, this time to the proof code check, implies that for every $x_0, x_1$,
there exists a measurement $\{H^{x_0, x_1}_w\}_w$ with outcomes in $\codes_{\ell_2}$ such that
\begin{equation}\label{eq:tweedledum}
M^{x_0, x_1, U}_a \otimes I_{\reg{Bob}} \approx_{\delta(\eps)} H^{x_0, x_1}_{[w|_{U} = a]} \otimes I_{\reg{Bob}}
\end{equation}
with respect to the distribution on~$\bx_0, \bx_1, \bU$.
Thus, by \Cref{fact:approx-delta-game-value}, we can assume that \Cref{eq:tweedledee,eq:tweedledum} hold with equality
with a loss of only $\delta(\eps)$ in the game value.
In addition, by \Cref{thm:naimark}, we can assume that the~$G$ and~$H$ measurements are all projective,
possibly replacing~$\psi$ with a different state.

\paragraph{Cross checks.}
Our next step is to apply the cross-checks. Passing these with probability $1-\delta(\eps)$ implies the bounds
\begin{equation}\label{eq:apply-the-cross-checks}
M^{x_0, x_1, T_0, T_1, U}_{a_0} \otimes I_{\reg{Bob}}\simeq_{\delta(\eps)} I_{\reg{Alice}} \otimes M^{x_0, T_0}_{a_0}
= I_{\reg{Alice}} \otimes G^{x_0}_{[w|_{T_0} = a_0]},
\end{equation}
\begin{equation}\label{eq:apply-the-cross-checks-2}
M^{x_0, x_1, T_0, T_1, U}_{a_1} \otimes I_{\reg{Bob}}\simeq_{\delta(\eps)} I_{\reg{Alice}} \otimes  M^{x_1, T_1}_{a_1}
= I_{\reg{Alice}} \otimes G^{x_1}_{[w|_{T_1} = a_1]},
\end{equation}
\begin{equation*}
M^{x_0, x_1, T_0, T_1, U}_{a_2} \otimes I_{\reg{Bob}}\simeq_{\delta(\eps)} I_{\reg{Alice}} \otimes  M^{x_0, x_1, U}_{a_2}
= I_{\reg{Alice}} \otimes H^{x_0, x_1}_{[w|_{U} = a_2]},
\end{equation*}
\begin{equation}\label{eq:Ms-self-consistency}
M^{x_0, x_1, T_0, T_1, U}_{a_0, a_1, a_2} \otimes I_{\reg{Bob}}\simeq_{\delta(\eps)} I_{\reg{Alice}} \otimes  M^{x_0, x_1, T_0, T_1, U}_{a_0, a_1, a_2}.
\end{equation}
At this point, we would like to apply \Cref{fact:low-degree-sandwich-on-steroids}.
To do so, we have to verify the distance property of our functions,
and this will follow from the fact that we augmented our index sets $\bI_0$, $\bI_1$, and $\bJ'$ with an additional uniformly random index.
To see this, consider two nonequal $w$ and $w'$ in $\codes_{\ell_1}$.
Then for them to agree on $\bT_0$, they must agree on $\bi_0$, and this happens only $\eta(\ell_1)$ fraction of the time.
The same holds for $\bU$, with the bound of $\eta(\ell_2)$.
As a result, \Cref{fact:low-degree-sandwich-on-steroids} implies the following:
consider the POVM measurement $\{\Lambda^{x_0, x_1}_{w_0, w_1, \pi}\}$
with outcomes~$w_0, w_1$ in $\codes_{\ell_1}$ and~$\pi$ in $\codes_{\ell_2}$
defined as
\begin{equation}\label{eq:dont-touch-me-bro}
\Lambda^{x_0, x_1}_{w_0, w_1, \pi}
:= G^{x_0}_{w_0} \cdot G^{x_1}_{w_1} \cdot H^{x_0, x_1}_{\pi} \cdot G^{x_1}_{w_1} \cdot G^{x_0}_{w_0}.
\end{equation}
Then
\begin{equation}\label{eq:gonna-use-this-once}
M^{x_0, x_1, T_0, T_1, U}_{a_0, a_1, a_2} \otimes I_{\reg{Bob}}\consistency_{\delta(\eps)}
I_{\reg{Alice}} \otimes \Lambda^{x_0, x_1}_{[w_0|_{T_0}, w_1|_{T_1}, \pi|_{U} = a_0, a_1, a_2]}.
\end{equation}
From this, \Cref{eq:Ms-self-consistency} implies
\begin{equation}\label{eq:gonna-use-this-once}
M^{x_0, x_1, T_0, T_1, U}_{a_0, a_1, a_2} \otimes I_{\reg{Bob}}\approx_{\delta(\eps)}
\Lambda^{x_0, x_1}_{[w_0|_{T_0}, w_1|_{T_1}, \pi|_{U} = a_0, a_1, a_2]} \otimes I_{\reg{Bob}}.
\end{equation}
Thus, by \Cref{fact:approx-delta-game-value}, we can assume that \Cref{eq:gonna-use-this-once} holds with equality by replacing~$M$ with~$G$,
incurring a loss of only $\delta(\eps)$ in the game value.
(Unlike before, here we do \emph{not} invoke \Cref{thm:naimark} on~$J^{x_0, x_1}$ to make it a projective measurement,
as that would likely change the structure in \Cref{eq:dont-touch-me-bro}, which we will need later.)

\paragraph{Verification.}
The strategy passes the verify check with probability $1-\delta(\eps)$.  
By \Cref{eq:gonna-use-this-once}
(which we now assume is equality),
this is the same probability as if we (i) sample~$\bx_0, \bx_1$,
(ii) use~$\Lambda$ to draw $\bw_0, \bw_1, \bpi$,
(iii) draw $\bI_0, \bI_1, \bJ$ conditioned on~$\bx_0, \bx_1$,
(iv) then draw $\bT_0, \bT_1, \bU$ conditioned on~$\bI_0, \bI_1, \bJ$,
(v) compute $\ba_0 = \bw_0|_{\bT_0}$, $\ba_1 = \bw_1|_{\bT_1}$, and $\ba_2 = \bw_2|_{\bU}$,
and (vi) give $\ba_0|_{\bI_0}, \ba_1|_{\bI_1}, \ba_2|_{\bJ}$ to $V_{\mathrm{PCPP}}$ and accept if it accepts.

Condition on a fixed choice of~$x_0, x_1$ and a draw for~$w_0, w_1, \pi$.
The PCPP verifier receives answers to its~$\bI_0$ and~$\bI_1$ queries based on~$w_0$ and~$w_1$, which are in~$\codes_{\ell_1}$.
In addition, although~$\pi$ is in~$\codes_{\ell_2}$ and may not correspond to the encoding of an actual proof string,
the verifier only queries it at points in the image of the embedding~$\mu_{\ell_2}$.
As a result, the answers $V_{\mathrm{PCPP}}$ receives to its~$J$ queries are consistent with \emph{some} fixed proof string.
Thus, by \Cref{prop:new-soundness}, since $1-\eta(k) \geq 2\gamma$ for all~$k$, if the probability
the verifier accepts is greater than~$s$, then there are strings $y_0, y_1 \in \{0, 1\}^{\ell_1}$
such that $w_0 = \codee_{\ell_1}(y_0)$, $w_1 = \codee_{\ell_1}(y_1)$
and $V(\mathsf{input}, x_0, x_1, y_0, y_1) = 1$.
Averaging over all~$\bx_0, \bx_1$ and~$\bw_0, \bw_1, \bpi$, we conclude that
\begin{equation}\label{eq:absolute-unit-of-a-strategy}
\Pr[V(\mathsf{input}, \bx_0, \bx_1, \coded_{\ell_1}(\bw_0), \coded_{\ell_1}(\bw_1)) = 1] \geq \frac{1-\delta(\eps)-s}{1-s} = 1-\delta(\eps).
\end{equation}
Recall that the decoding map is one-to-one except on those strings not in the range of the encoding map, which it maps to~$\bot$ instead.
As we can assume that the verifier~$V$ always rejects when it receives~$\bot$ for an answer,
this tells us that $\coded_{\ell_1}(\bw_0), \coded_{\ell_1}(\bw_1) \neq \bot$ with probability at least $1-\delta(\eps)$.

\paragraph{Wrapping it up.}
Now we give a strategy for causing the verifier~$V$ to accept with high probability on~$\mathsf{input}$.
It uses state~$\psi$, and given question~$x$ it applies the measurement $\{A^x_a\}_a$ defined as
\begin{equation*}
A^x_a := G^x_{[\coded_{\ell_1}(w) = a]}.
\end{equation*}
Consider the verifier~$V'$ which samples $(\bx_0, \bx_1)$, gives them to Alice and Bob, receives~$\bw_0, \bw_1$,
and accepts if $V(\mathsf{input}, \bx_0, \bx_1, \coded_{\ell_1}(\bw_0), \coded_{\ell_1}(\bw_1)) = 1$.
Then~$V$ accepts on strategy~$A$ with the same probability that~$V'$ accepts on strategy~$G$.
In other words, if we define $S(x_0, x_1)$ to be the set of $(w_0, w_1)$
such that
\begin{equation*}
V(\mathsf{input}, x_0, x_1, \coded_{\ell_1}(w_0), \coded_{\ell_1}(w_1))=1,
\end{equation*}
then the probability that~$V'$ accepts on strategy~$G$ is
\begin{equation}\label{eq:its-the-final-countdown}
\E_{(\bx_0, \bx_1)} \sum_{w_0, w_1 \in S(\bx_0, \bx_1)} \bra{\psi} G^{x_0}_{w_0} \otimes G^{x_1}_{w_1}\ket{\psi}.
\end{equation}

To show this is large, we begin by showing that the~$G$'s commute with each other.
To see this, note that \Cref{eq:apply-the-cross-checks,eq:apply-the-cross-checks-2} implies that for a fixed $c \in \{0, 1\}$,
\begin{equation*}
\Lambda^{x_0, x_1}_{[w_c|_{T_c} = a_c]} \otimes I_{\reg{Bob}}\simeq_{\delta(\eps)}
 I_{\reg{Alice}} \otimes G^{x_c}_{[w_c'|_{T_c} = a_c]}.
\end{equation*}
However, by~the distance properties of our code and the fact that $\bT_c$ contains a uniformly random index, this implies that
\begin{equation}\label{eq:dont-copy-that-floppy}
\Lambda^{x_0, x_1}_{w_c} \otimes I_{\reg{Bob}}\simeq_{\delta(\eps)}
 I_{\reg{Alice}} \otimes G^{x_c}_{w_c}.
\end{equation}
As a result,
\begin{align*}
G^{x_0}_{w_0} \cdot G^{x_1}_{w_1} \otimes I_{\reg{Bob}}
&\approx_{\delta(\eps)} G^{x_0}_{w_0} \otimes \Lambda^{x_0, x_1}_{w_1}\\
&\approx_{\delta(\eps)} I_{\reg{Alice}} \otimes \Lambda^{x_0, x_1}_{w_1}\cdot \Lambda^{x_0, x_1}_{w_0}\\
&\approx_{\delta(\eps)} I_{\reg{Alice}} \otimes \Lambda^{x_0, x_1}_{w_0}\cdot \Lambda^{x_0, x_1}_{w_1}\\
&\approx_{\delta(\eps)} G^{x_1}_{w_1} \otimes \Lambda^{x_0, x_1}_{w_0}\\
&\approx_{\delta(\eps)}  G^{x_1}_{w_1} \cdot G^{x_0}_{w_0} \otimes I_{\reg{Bob}}.
\end{align*}
A similar argument as the one establishing~\eqref{eq:dont-copy-that-floppy} implies that
\begin{equation*}
G^{x_c}_{w_c} \otimes I_{\reg{Bob}}\simeq_{\delta(\eps)}
 I_{\reg{Alice}} \otimes G^{x_c}_{w_c}.
\end{equation*}
Thus,
\begin{align*}
G^{x_0}_{w_0} \otimes G^{x_1}_{w_1}
&= G^{x_0}_{w_0} \cdot G^{x_0}_{w_0} \otimes G^{x_1}_{w_1}\\
&\approx_{\delta(\eps)} G^{x_0}_{w_0} \cdot G^{x_0}_{w_0} \cdot G^{x_1}_{w_1} \otimes I_{\reg{Bob}}\\
&\approx_{\delta(\eps)} G^{x_0}_{w_0} \cdot G^{x_1}_{w_1} \cdot G^{x_0}_{w_0} \otimes I_{\reg{Bob}}\\
&= \Lambda^{x_0, x_1}_{w_0, w_1} \otimes I_{\reg{Bob}}.
\end{align*}
As a result, by \Cref{fact:approx-delta-generalized-game-value},  \Cref{eq:its-the-final-countdown} is at least $1-\delta(\eps)$ by \Cref{eq:absolute-unit-of-a-strategy}.
This concludes the proof of the theorem.
\end{proof}

\subsection{Applying the answer reduction protocol}

In this section, we instantiate \Cref{thm:answer-reduction} with the low-degree code
and then apply it to our $\neexp$ protocol.

\begin{theorem}\label{thm:applied-with-low-degree-code}
Let $V = (\mathrm{Alg}_{\mathrm{Q}}, \mathrm{Alg}_{\mathrm{A}})$ be an $\MIP^*$ verifier for a language~$L$.
		Write $L_{\mathrm{A}}$ for the language decided by $\mathrm{Alg}_{\mathrm{A}}$.
		Suppose on inputs of size~$n$, the verifier~$V$ has question length~$\ell_{V,\mathrm{Q}}(n)$, answer length~$\ell_{V,\mathrm{A}}(n)$,
		question time~$t_{V,\mathrm{Q}}(n)$, and answer time~$t_{V,\mathrm{A}}(n)$.
		Then there exists another $\MIP^*$ verifier $V_{\mathrm{ans}}$ for~$L$ with the following parameters.
\begin{align*}
\qlength{V_{\mathrm{ans}}} &= O(\ell_{V, \mathrm{Q}}(n) + \log(\ell_{V,\mathrm{A}}(n))+ \log(t_{V, \mathrm{A}}(n))),\\
\alength{V_{\mathrm{ans}}} &= O(\mathrm{poly log}(\ell_{V, \mathrm{A}}(n)) + \mathrm{poly log}(t_{V, \mathrm{A}}(n))),\\
\qtime{V_{\mathrm{ans}}} &= O\left(t_{V,\mathrm{Q}}(n)) + \poly(n + \ell_{V, \mathrm{Q}}(n), \log(\ell_{V, \mathrm{A}}(n)), \log(t_{V, \mathrm{A}}(n))\right),\\
\atime{V_{\mathrm{ans}}} &= \poly(n + \ell_{V, \mathrm{Q}}(n), \log(\ell_{V, \mathrm{A}}(n)), \log(t_{V, \mathrm{A}}(n))).
\end{align*}
\end{theorem}
\begin{proof}
We instantiate the low-degree code in \Cref{fact:low-degree-code-is-efficient}.
It gives an error correcting code with parameters $(n, \poly(n), \mathrm{poly log}(n), \mathrm{poly log}(n)^{-1},\poly(n), \mathrm{polylog}(n))$
and a $c$-subset test~$\game_k$ with robustness $\chi_k(\eps) = \poly(\eps, \log(k)^{-1})$ such that
\begin{equation*}
\qtime{\game_k} = \poly(\log k, c),
\quad
\atime{\game_k} = \poly(\log(k)^c),
\end{equation*}
\begin{equation*}
\qlength{\game_k} = O(c \log k),
\quad
\alength{\game_k} = O(\log(k)^{2c}).
\end{equation*}
We then apply \Cref{thm:answer-reduction} with $s, \gamma = \tfrac{1}{10}$.
At this point, the theorem follows immediately,
but as deriving it can be cumbersome, we fill in the details.

By construction, $t_{\mathrm{Dec}}(n) = \poly(n)$.
As a result,
\begin{equation*}
t_{\mathrm{compose}}(n)
= t_{V, \mathrm{A}}(n) + t_{\mathrm{Dec}}(\ell_{V,\mathrm{A}}(n))
= t_{V, \mathrm{A}}(n) + \poly(\ell_{V,\mathrm{A}}(n)).
\end{equation*}
Thus,
\begin{equation*}
\ell_\pi(n)
= t_{\mathrm{compose}}(n) \cdot \mathrm{poly log}(t_{\mathrm{compose}}(n))
= \poly(t_{V, \mathrm{A}}(n), \ell_{V,\mathrm{A}}(n)).
\end{equation*}
Now, $m(n) = \poly(n)$. Thus,
\begin{align*}
t_{\mathrm{PCPP}}(n)
&= \poly(n + \ell_{V, \mathrm{Q}}(n), \log(m(\ell_{V, \mathrm{A}}(n))), \log(t_{\mathrm{compose}}(n)))\\
&= \poly(n + \ell_{V, \mathrm{Q}}(n), \log(\ell_{V, \mathrm{A}}(n)), \log(t_{V, \mathrm{A}}(n))).
\end{align*}
Furthermore, $q(n) = \mathrm{polylog}(n)$ and $t_{\mathrm{Emb}}(n) = \mathrm{poly log}(n)$.
As a result,
\begin{align*}
\log(m(\ell_\pi(n))) &= O(\log(\ell_{V,\mathrm{A}}(n))+ \log(t_{V, \mathrm{A}}(n))),\\
\log(q(\ell_\pi(n))) &= O(\log \log(\ell_{V,\mathrm{A}}(n))+ \log \log(t_{V, \mathrm{A}}(n))),\\
t_{\mathrm{Emb}}(\ell_\pi(n)) &= \poly(\log(\ell_{V, \mathrm{A}}(n)), \log(t_{V, \mathrm{A}}(n))).
\end{align*}
The theorem now follows from applying these bounds to \Cref{thm:answer-reduction}.
\end{proof}

Crucially, although polynomial factors of $t_{V,Q}(n)$ and $\ell_{V, Q}(n)$ appear in \Cref{thm:applied-with-low-degree-code},
only the \emph{logarithms} of $t_{V,A}(n)$ and $\ell_{V, A}(n)$ appear in this theorem.
As a result, if we apply this to \Cref{cor:succinct-sat-protocol-with-big-answer-size}, we arrive at our main result.

\begin{theorem}
  \label{thm:main}
There is an $\MIP^*$ verifier~$\game$ for $\succinctsquared$ with parameters
\begin{equation*}
\qlength{\game} =O(n), \quad
\alength{\game} = \poly(n),
\end{equation*}
\begin{equation*}
\qtime{\game} = \poly(n), \quad
\atime{\game} = \poly(n).
\end{equation*}
\end{theorem}

%%% Local Variables:
%%% mode: latex
%%% TeX-master: "main"
%%% End:
%%%---------- close: answer-reduction






